\documentclass[10pt]{amsart}
\usepackage[a4paper,margin=23mm]{geometry}
\usepackage{amsmath,amssymb,amsthm,mathtools,hyperref,graphicx,xspace,letltxmacro,xpatch}
\usepackage[english]{babel}
\newtheorem{theo}{Theorem}[section]
\newtheorem{lemm}[theo]{Lemma}\newtheorem{prop}[theo]{Proposition}\newtheorem{coro}[theo]{Corollary}\newtheorem{conj}[theo]{Conjecture}
\theoremstyle{definition}\newtheorem{defi}[theo]{Definition}\newtheorem{rema}[theo]{Remark}\newtheorem{exam}[theo]{Example}\newtheorem{exem}[theo]{Example}\newtheorem{nota}[theo]{Notation}
\newcommand{\ie}{i.e.}\newcommand{\eg}{e.g.}\newcommand{\resp}{resp.}\newcommand{\skpt}{\par\smallskip\noindent}
\emergencystretch=3em\pagestyle{plain}

%%\usepackage{amssymb}
%%\usepackage{amsfonts}
%\usepackage{amsthm}
%\usepackage{amsmath}
%%\usepackage{mathtools}
%%\usepackage{graphicx}
\usepackage{mathtools}

\makeatletter
\def\@altabstractlanguage{}
\makeatother

\newcommand*{\dd}{\mathrm{d}}
\newcommand*{\dx}{\dd x}
\newcommand*{\dy}{\dd y}
\newcommand*{\dz}{\dd z}
\newcommand*{\du}{\dd u}
\newcommand*{\dv}{\dd v}
\newcommand{\sdots}{{\mk.\mK.\mK.\mK}}
\newcommand{\scdots}{{\mk{\cdot}\mK{\cdot}\mK{\cdot}\mk}}

\newcommand*{\myangle}{\mathord{\langle}} %~ \langle est une pseudo parenthèse mais l'auteur l'utilise comme un symbole ordinaire. Il faut donc changer son type pour corriger les problèmes d'espace avec les +, -, etc.

\newcommand*{\I}{\textsl{I}}
\newcommand*{\II}{\textsl{II}} % pour mieux gérer les espaces entre les I

%\usepackage{amscd}


\DeclareMathOperator{\airy}{Airy} 
\DeclareMathOperator{\hit}{hit} 
\DeclareMathOperator{\shk}{shk} 
\DeclareMathOperator{\tw}{TW} 
\DeclareMathOperator{\err}{Err} 
\DeclareMathOperator{\hypo}{hypo} 
\DeclareMathOperator{\tr}{tr} 
\DeclareMathOperator{\hs}{HS} 
\DeclareMathOperator{\op}{op} 
\DeclareMathOperator{\Ai}{Ai}
\DeclareMathOperator{\argl}{arg} 
\DeclareMathOperator{\const}{const} 


\newcommand{\E}[1]{\mathbf{E}\left[#1\right]}
\newcommand{\pr}[1]{\mathbf{Pr}\left[#1\right]}
\newcommand{\ind}[1]{\mathbf{1}_{\{ #1 \}}}
\newcommand{\eps}{\varepsilon}
\newcommand{\Z}{\mathbb{Z}}
\newcommand{\R}{\mathbb{R}}
%%\newcommand{\Ai}{\mathrm{Ai}}
\newcommand{\A}{\mathcal{A}}
\newcommand{\h}{\mathfrak{h}}

%%%%%%%%%%%%%%%%%%%%%%%%%%%%%%%%%%%%%%%%%%%%%%%%%%%%%%%%%%%%%%%%%%%%
\providecommand{\noopsort}[1]{}
\graphicspath{{./figures/}}

\newcommand*{\mk}{\mkern -1mu}
\newcommand*{\Mk}{\mkern -2mu}
\newcommand*{\mK}{\mkern 1mu}
\newcommand*{\MK}{\mkern 2mu}

\hypersetup{urlcolor=purple, linkcolor=blue, citecolor=red}

\newcommand*{\romanenumi}{\renewcommand*{\theenumi}{\roman{enumi}}}
\newcommand*{\Romanenumi}{\renewcommand*{\theenumi}{\Roman{enumi}}}
\newcommand*{\alphenumi}{\renewcommand*{\theenumi}{\alph{enumi}}}
\newcommand*{\Alphenumi}{\renewcommand*{\theenumi}{\Alph{enumi}}}
%%%%%%%%%%%%%%%%%%%%%%%%%%%%%%%%%%%%%%%%%%%%%%%%%%%%%%%%%%%%%%%%%%%%

\newtheorem*{rema*}{Remark}\NewDocumentEnvironment{enonce*}{ O{} m O{} }{\par\smallskip\noindent\textbf{#2\if\relax\detokenize{#3}\relax\else\ (#3)\fi.}\ }{\par\smallskip}
